% !TeX document-id = {21d95b9e-6149-4bf7-a0af-622fadedeb8b}
% !TeX TXS-program:compile = txs:///pdflatex/[--shell-escape]
\documentclass[spanish,addpoints,answers,a4paper]{exam}
\usepackage[T1]{fontenc}
\usepackage[spanish]{babel}
\usepackage{libertine}
\usepackage[tmargin=2cm,rmargin=2cm,lmargin=2cm]{geometry}
\usepackage{amsmath,bm}
\usepackage[shortlabels]{enumitem}
\usepackage{siunitx}
\usepackage{multicol}

\usepackage{graphicx}
\graphicspath{{../../../images/}}
\pagestyle{plain}

\usepackage{minted}
\usepackage{hyperref}
\renewcommand{\solutiontitle}{\noindent\textbf{Solución}\par\noindent}
\renewcommand\listingscaption{Listado}
\newcommand{\unmarkedfntext}[1]{%
	\begingroup
	\renewcommand\thefootnote{}\footnote{#1}%
	\addtocounter{footnote}{-1}%
	\endgroup
}
\spanishdatedel
\decimalpoint

\def\LOGOUNI{%
	\begin{picture}(0,0)\unitlength=1cm
	\put (-13.5,-2.5) {\includegraphics[width=2.8cm]{logouni}}
	\end{picture}
}

\def\LOGOCTIC{%
	\begin{picture}(0,0)\unitlength=1cm
	\put (-8,-1.2) {\includegraphics[height=1cm]{logocticblack}}
	\end{picture}
}

\newcommand{\mychar}{%
	\begingroup\normalfont
	\includegraphics[height=\fontcharht\font`\B]{Octocat.png}%
	\endgroup
}
\begin{document}

\begin{center}
	\sffamily\bfseries\scshape
	{\Large UNIVERSIDAD NACIONAL DE INGENIERÍA}\LOGOUNI\\
	Centro de Tecnologías de la Información y Comunicaciones\LOGOCTIC\\
\end{center}

\vspace{.8cm}

\begin{center}\sffamily\bfseries\large
	C/C++ PROGRAMMING LANGUAGE\\
	TEMA $\bm{4}$: OPERATORS, EXPRESSIONS, AND STATEMENTS 
\end{center}

\vspace{.5cm}
\noindent
\makebox[\textwidth]{Nombres y apellidos:\enspace\hrulefill}
\vspace{0.2in}
\makebox[\textwidth]{Nombres y apellidos del instructor:\enspace MSc. César Manuel Sebastián Díez Chirinos.\hfill}

\begin{questions}

\question Asuma que todas las variables son del tipo \mintinline{c}{int}. Encuentre el valor de cada una de las siguientes variables:

\begin{parts}
\part \mintinline{c}{x = (2 + 3) * 6;}
\part \mintinline{c}{x = (12 + 6)/2*3;}
\part \mintinline{c}{y = x = (2 + 3)/4;}
\part \mintinline{c}{y = 3 + 2*(x = 7/2);}
\end{parts}

\begin{solution}
Si tenemos el siguiente listing~\ref{lst:4.1}, obtendremos para los valores de \mintinline{c}{x} e \mintinline{c}{y}:

\begin{listing}[H]
	\footnotesize
	\inputminted{c}{exercise4_1.c}
	\caption{Programa \texttt{exercise4\_1.c}.}
	\label{lst:4.1}
\end{listing}


\begin{multicols}{2}
\begin{enumerate}[(a)]

\item Imprimirá:

\begin{minted}{bash}
pc@CTIC:~$ x: 30
\end{minted}

\item Imprimirá:

\begin{minted}{bash}
pc@CTIC:~$ x: 27
\end{minted}

\item Imprimirá:

\begin{minted}{bash}
pc@CTIC:~$ x: 1    y: 1
\end{minted}

\item Imprimirá:

\begin{minted}{bash}
pc@CTIC:~$ x: 3    y: 9
\end{minted}

\end{enumerate}
\end{multicols}

\end{solution}

\question Sospechas que hay algunos errores en el siguiente programa. ¿Puede encontrarlos?

\begin{minted}[linenos]{c}
int main(void)
{
	int=1,
	float n;
	printf("watch out! Here come a bunch of fractions!\n");
	while(i<30)
		n=1/i;
		printf(" %f", n);
		
printf("That's all, folks!\n");
return;
}
\end{minted}

\begin{solution}
\begin{enumerate}[label={Línea \arabic*:},leftmargin=1.2cm]\setcounter{enumi}{-1}
	\item Debe ir la cabecera \mintinline{c}{#include <stdio.h>}.\setcounter{enumi}{2}
	\item Debe acabar en punto y coma, no en coma.\setcounter{enumi}{5}
	\item La instrucción \mintinline{c}{while} establece un ciclo infinito porque el valor de \mintinline{c}{i} permanece \mintinline{c}{1}	y siempre es menos de \mintinline{c}{30}. Presumiblemente, se quiso escribir \mintinline{c}{while (i++ <30)}.
	\item[Línea 6--8] La sangría implica que queríamos las líneas 7 y 8 para formar un bloque, pero la falta de llaves significa que el ciclo while incluye solo la línea 7. Deben agregarse llaves.\setcounter{enumi}{6}
	\item Como \mintinline{c}{1} e \mintinline{c}{i} son ambos enteros, el resultado de la división será \mintinline{c}{1} cuando \mintinline{c}{i} es \mintinline{c}{1} y \mintinline{c}{0} para todos los valores más grandes. La sentencia \mintinline{c}{n = 1.0 / i;} causaría que \mintinline{c}{n} se convierta en flotante.
	\item Omitimos un carácter de nueva línea \mintinline{c}{\n} en la declaración de control. Esto causa los números sean impresos en una sola línea.\setcounter{enumi}{9}
	\item Debería estar \mintinline{c}{return 0;}.
\end{enumerate}

Aquí hay una posible versión correcta:

\begin{listing}[H]
	\footnotesize
	\inputminted{c}{exercise4_2.c}
	\caption{Programa \texttt{exercise4\_2.c}.}
	\label{lst:4.2}
\end{listing}
\end{solution}

\question Hacer un \mintinline{c}{min_sec} interactivo no es fácil. ¿Cómo se puede mejorar?

\begin{minted}{c}
#include<stdio.h>
#define S_TO_M 60
main()
{
int sec, min, left;
printf("This program convierte segundos a minutos y");
printf("segundos.\n");
printf("Solo debe ingresar los segundos.\n");
printf("Ingrese 0 para finalizar el programa.\n");
while(sec>0){
scanf("%d", &sec);
min=sec/S_TO_M;
left=sec%S_TO_M;
printf("%d sec is %d min, %d sec.\n", sec, min, left);
printf("Next input?\n");
}
printf("See you!\n");
}
\end{minted}

\begin{solution}

El principal problema radica en la relación entre la declaración de prueba (¿es \mintinline{c}{sec} mayor que $0$?) Y la instrucción \mintinline{c}{scanf()} que obtiene el valor de \mintinline{c}{sec}. En particular, la primera vez que se realiza la prueba, el programa no ha tenido la oportunidad de obtener un valor por segundo, y la comparación se realizará con algún valor basura que se encuentre en esa ubicación de memoria; puede o no ser mayor que $0$. Una solución, aunque poco elegante, es inicializar \mintinline{c}{sec} a, por ejemplo, \mintinline{c}{1} para que la prueba se pase la primera vez. Esto descubre un segundo problema. Cuando finalmente escribes \mintinline{c}{0} para detener el programa, \mintinline{c}{sec} no
verifique hasta que termine el ciclo y se impriman los resultados de $0$ segundos. Lo que realmente quiere es tener una instrucción \mintinline{c}{scanf()} justo antes de realizar la prueba. Puede lograr eso alterando la parte central del programa para leer de esta manera:

\begin{listing}[H]
	\footnotesize
	\inputminted{c}{exercise4_3.c}
	\caption{Programa \texttt{exercise4\_3.c}.}
	\label{lst:4.3}
\end{listing}

La primera vez, se utiliza el \mintinline{c}{scanf()} fuera del bucle. A partir de entonces, se usa el \mintinline{c}{scanf()} al final del ciclo (y, por tanto, justo antes de que el ciclo comience de nuevo).
\end{solution}

\question Escriba un programa que pregunte por un entero, y que imprima todos los enteros desde este número hasta $10$ más de este. (Si fuera $5$, sería desde $5$ hasta $15$). Asegúrese de separar cada valor de salida por un espacio o tabulador o una nueva línea.

\begin{solution}

\begin{listing}[H]
	\footnotesize
	\inputminted{c}{exercise4_4.c}
	\caption{Programa \texttt{exercise4\_4.c}.}
	\label{lst:4.4}
\end{listing}
\end{solution}

\question Escriba un programa que solicite un decimal e imprima su cubo.

\begin{solution}

\begin{listing}[H]
	\footnotesize
	\inputminted{c}{exercise4_5.c}
	\caption{Programa \texttt{exercise4\_5.c}.}
	\label{lst:4.5}
\end{listing}
\end{solution}

\question Use un \mintinline{c}{while} loop para convertir el tiempo en minutos a el tiempo en horas y minutos.

\begin{solution}
	
\begin{listing}[H]
	\footnotesize
	\inputminted{c}{exercise4_6.c}
	\caption{Programa \texttt{exercise4\_6.c}.}
	\label{lst:4.6}
\end{listing}
\end{solution}

\question ¿Qué imprimirá este programa?

\begin{minted}{c}
#include<stdio.h>
#define FORMAT "%s is a string\n"
main()
{
int num=0;
printf(FORMAT, FORMAT);
printf("%d\n", ++num);
printf("%d\n", num++);
printf("%d\n", num--);
printf("%d\n", num);
}
\end{minted}

\begin{solution}

\begin{listing}[H]
	\footnotesize
	\inputminted{c}{exercise4_7.c}
	\caption{Programa \texttt{exercise4\_7.c}.}
	\label{lst:4.7}
\end{listing}


Luego de ejecutar el programa, se obtiene en consola el resultado:

\begin{minted}{bash}
pc@CTIC:~$ x: gcc exercise4_7.c && ./a.out
%s is a string
is a string
1
1
2
1
\end{minted}

\end{solution}

\question Cambie el programa \texttt{addemup.c} para calcular cuánto dinero ganaría en $20$ días, si recibe $1\$$ el primer día, $2\$$ el segundo, $3\$$ el tercero y así.

\begin{solution}

\begin{listing}[H]
	\footnotesize
	\inputminted{c}{exercise4_8.c}
	\caption{Programa \texttt{exercise4\_8.c}.}
	\label{lst:4.8}
\end{listing}
\end{solution}

\question Escriba un programa que convierta sus días en semanas y días.

\begin{solution}

\begin{listing}[H]
	\footnotesize
	\inputminted{c}{exercise4_9.c}
	\caption{Programa \texttt{exercise4\_9.c}.}
	\label{lst:4.9}
\end{listing}
\end{solution}

\question Construya sentencia que hagan lo siguiente:

\begin{enumerate}[a.]
	\item Incremente la variable \mintinline{c}{x} por 10.
	\item Incremente la variable \mintinline{c}{x} por 1.
	\item Asigne dos veces la suma de \mintinline{c}{a} y \mintinline{c}{b} a \mintinline{c}{c}.
	\item Asigne \mintinline{c}{a} más dos veces \mintinline{c}{b} a \mintinline{c}{c}.
\end{enumerate}

\begin{solution}

\begin{listing}[H]
	\footnotesize
	\inputminted{c}{exercise4_10.c}
	\caption{Programa \texttt{exercise4\_10.c}.}
	\label{lst:4.10}
\end{listing}
\end{solution}

\end{questions}

\begin{flushright}\bfseries
Centro de Tecnologías de la Información y Comunicaciones (CTIC)\\[2mm]
\today\unmarkedfntext{Código disponible en \href{https://github.com/carlosal1015/C-Programming}{\mychar{}}.}
\end{flushright}

\end{document}